\subsubsection{雷达站是什么}

\textbf{雷达}方向：利用场地高处视角与算力，识别全场目标等信息，经裁判系统向队友共享。
它与导航章中的\textbf{激光雷达（LiDAR）}不是同一概念。

与自瞄的核心差异：
\begin{itemize}
    \item 自瞄：单车、近距、云台闭环、持续跟瞄；
    \item 雷达：全场、多目标、俯视几何、\textbf{广播}为主，不直接拧本车云台。
\end{itemize}

\begin{example}[区分站与传感器]
写下“高台相机识别目标并共享坐标”和“车载 LiDAR 量测墙面帮助定位”，圈出最终输出：前者是全场目标信息，后者服务本车位姿。检查：名称含“雷达”并不能说明承担同一种比赛任务。
\end{example}

\begin{exercise}
\label{v3-ex-radar-1}
一套高台视觉系统没有 LiDAR，还可以属于雷达方向吗？说明依据。
\end{exercise}

\subsubsection{典型软件架构（概述）}

可把软件拆成以下教学架构；其中无线电与反制均为可选示意，不代表每个赛季都有相同任务或权限：
\begin{enumerate}
    \item \textbf{地面视觉}：工业相机 → 检测/跟踪 → 场地图坐标 → 融合 → 裁判串口；
    \item \textbf{无线电（可选）}：软件定义无线电（Software Defined Radio，SDR）按适用协议接收与解析信息，经本机 TCP 交给视觉；
    \item \textbf{反制（可选）}：长焦检测与云台，和地面端用独立串口同步次数。
\end{enumerate}

通信模式见卷二「多机通信」；SDR 与解析见卷二「SDR 与电磁波解析」。本卷只要求能画数据流并说出职责边界。

数据流可记为：工业相机经识别跟踪得到场地图坐标；可选的 SDR 解码坐标按新鲜度与视觉融合后，经裁判系统广播；反制计数为可选旁路。

\begin{figure}[H]
\centering
\resizebox{0.9\textwidth}{!}{%
\begin{tikzpicture}[box/.style={draw=cyan!60!black,rounded corners,align=center,
  minimum height=11mm,text width=27mm,fill=cyan!5},>=stealth]
\node[box] (vision) at (0,1) {视觉检测与映射\\场地系坐标};
\node[box,dashed] (radio) at (0,-1) {可选无线电解析\\有效坐标};
\node[box] (fusion) at (4.2,0) {身份、质量与时效\\检查后选源};
\node[box] (send) at (8.4,0) {协议封装与广播\\坐标或无效状态};
\draw[->,thick] (vision.east) -- (fusion.west);
\draw[->,thick,dashed] (radio.east) -- (fusion.west);
\draw[->,thick] (fusion) -- (send);
\end{tikzpicture}}
\caption{雷达站坐标主路；需按序处理的事件不与最新坐标共用覆盖队列}
\label{fig:v3-radar-chain}
\end{figure}

\begin{example}[把旁路从主路分开]
画出视觉进程 → 融合 → 广播三框，再从可选无线电进程画箭头接到融合。用手遮住无线电框，检查主路是否仍能输入视觉坐标；再给非坐标事件单画一条有序处理支路。检查：可选模块故障不应在设计上阻塞仍有效的坐标主路。
\end{example}

\begin{exercise}
\label{v3-ex-radar-2}
为何坐标可以只保留最新值，反制计数或任务事件却不能随意覆盖？
\end{exercise}

\subsubsection{俯视定位 vs 车载 PnP}

雷达可用场地图关键点与单应变换（Homography）或区域映射，把适用平面上的图像点映射到场地坐标；有高度差时需分区域、补充高度信息或使用三维方法验证，不能把任意车顶像素都当作地面点。
自瞄常用装甲板模型 + PnP 求相对位姿。二者标定物、误差来源与调参方式都不同，不要混用术语。

检测上常见级联：车辆检测 → 装甲/灯条 → 数字分类；再接多目标跟踪。细节回卷二 B/C。

\begin{example}[验证平面假设]
在教学场地图上标出四个任意三点不共线的地面角点，声明它们是图像与场地的对应点。再选一个未参与标定的地面点检查映射误差。把验证点移到高台上后，标注“高度不在同一平面，原映射不保证有效”。检查：映射适用范围必须写在配置旁。
\end{example}

\begin{exercise}
\label{v3-ex-radar-3}
直接把车顶检测框中心用地面单应映射到场地，为什么可能产生偏差？
\end{exercise}

\subsubsection{融合与时效}

若存在无线电坐标，可采用一种教学策略：先检查身份、坐标系、质量与时效，再让\textbf{合格且新鲜的无线电优先}；否则回退有效视觉或显式标记无效。这不是所有系统通用的优先级。
非坐标事件可走独立处理通道，避免阻塞坐标。坐标可采用只保留最新值（latest-value）的方式；计数增量、事件或需确认的消息不能直接套用覆盖策略。

对队友有用的是「对的坐标系 + 可接受延迟」；错系、堵包、过期坐标比「偶尔缺一帧」更具破坏性。

\begin{example}[按有效性与时效选源]
教学策略：同一目标、同一坐标系且质量合格时，无线电年龄不超过 100 ms 则优先；否则选择年龄不超过 150 ms 的视觉；二者均不合格就输出无效。当前时刻 1000 ms，无线电采样于 850 ms，视觉采样于 950 ms，年龄分别为 150、50 ms，应选视觉。这些阈值仅用于演练。
\end{example}

\begin{note}
本章不提供当年比赛的功能授权、频率或协议常量。无线电、反制与裁判广播须由负责人核对参赛赛季官方规则与协议后采用。平面映射前提见 OpenCV 官方说明：\url{https://docs.opencv.org/4.x/d9/dab/tutorial_homography.html}。
\end{note}

\begin{conclusion}
\begin{itemize}
  \item 先确认同一目标与坐标，再检查质量与年龄，最后选源。
  \item 接收时刻不等于采样时刻。
  \item 两源均无效时明确报告失效。
\end{itemize}
\end{conclusion}

\begin{exercise}
\label{v3-ex-radar-4}
沿用示例策略，当前 1000 ms，无线电采样 930 ms 但校验失败，视觉采样 960 ms 且有效，选谁？若视觉也过期呢？
\end{exercise}

\subsubsection{技能映射}

\begin{center}
\begin{tabular}{@{}p{0.30\textwidth}p{0.62\textwidth}@{}}
\toprule[1.5pt]
能力 & 卷二 \\
\midrule[1pt]
节点与工程 & ROS2 \\
TCP/串口/裁判帧 & 多机通信 \\
SDR/IQ/状态机 & SDR 与电磁波解析 \\
检测与跟踪 & OpenCV、深度学习 \\
相机与标定 & 相机成像 \\
部署加速 & CUDA \\
\bottomrule[1.5pt]
\end{tabular}
\end{center}

\begin{example}[沿故障查专项]
教学问题是“识别框稳定，队友收到的位置总是很旧”。先沿发送队列、融合取值、接收消费画检查点，再到卷二多机通信查最新值与故障隔离。检查：此时优先补通信知识，比立即换检测模型更直接。
\end{example}

\begin{exercise}
\label{v3-ex-radar-5}
若画面上目标清楚但数字类别经常错，应先查卷二哪些章？若坐标正确但串口帧被拒收呢？
\end{exercise}
